\section{Socio-Technical Implications: The Industry Structure of Intelligent Computing}
\label{sec:sociotechnical}

The ICA framework describes a layered architecture for model-native computing systems. Yet architecture does not exist in a vacuum: it both shapes and is shaped by the economic structure of the industry that builds it. The semiconductor and personal-computer industries underwent a well-documented disaggregation over four decades, evolving from vertically integrated giants into layered ecosystems of specialized firms. The emerging AI industry appears to be following a parallel trajectory, with one critical difference: model weights are software, infinitely copyable at near-zero marginal cost, which means the physical moat that once protected chipmakers has no direct analog. This section traces the structural dynamics that will determine whether intelligent computing settles into a disaggregated stack or remains dominated by a small number of vertically integrated players.

\subsection{From IDM to Fabless: The Semiconductor Analogy and AI Stack Disaggregation}
\label{sec:idm-fabless}

The semiconductor industry evolved through four structural stages. Before the 1980s, firms such as Intel, AMD, and Texas Instruments operated as Integrated Device Manufacturers (IDMs), controlling every stage from design through fabrication to sales~\cite{generativevalue2024semiconductor}. TSMC's founding in 1987 introduced the pure-play foundry model, separating fabrication from design~\cite{quartr2024tsmc}. The fabless era followed: NVIDIA and Qualcomm designed chips without owning fabs. ARM then added a fourth layer by licensing processor IP without manufacturing at all~\cite{eetimes2025asic}. Each stage lowered entry barriers and created specialized value-chain layers.

The AI industry today resembles the IDM phase. NVIDIA dominates as a vertically integrated platform, owning both the hardware (GPUs) and the software stack (CUDA), much as IBM controlled mainframe hardware, operating systems, and applications before the PC era's disaggregation. Ben Evans frames the central structural question: will AI follow the PC-era pattern of disaggregation into separate chip, platform/OS, model, and application layers, or will it remain vertically integrated under dominant players?~\cite{benevans2025fullstack}

The first wave of disaggregation is already visible. Cloud giants building custom silicon, such as Google TPUs, Amazon Trainium, and Microsoft Maia, signal the separation of hardware provision from the software and model layers, analogous to how the foundry model separated fabrication from design~\cite{industryanalysis2026aivaluechain}. Stanford HAI's 2025 data showing a 280-fold drop in reasoning costs, combined with surging small-model performance, creates commoditization pressure on foundation model providers~\cite{stanfordhai2025index}, accelerating the structural parallel to how commodity x86 CPUs enabled the PC disaggregation. Gartner's forecast of major consolidation in agentic AI due to oversupply mirrors the semiconductor industry's consolidation cycles~\cite{gartner2026trends}, suggesting that only a few foundation model providers will survive to become the ``Intel/AMD'' of AI.

\subsection{Model Weights as Software: The Copyability Problem and the Closed-Source Strategic Response}
\label{sec:copyability}

The semiconductor analogy, however, breaks down at one crucial point. A CPU's intellectual property is embedded in its physical manufacturing process: reverse-engineering a chip from its external behavior is economically infeasible. Model weights, by contrast, are fundamentally software. They can be copied infinitely at near-zero marginal cost, and they can be extracted from API outputs through model distillation attacks. This asymmetry is the primary structural reason that top-tier models (GPT-4, Claude, Gemini) remain closed-source.

The closed-source strategy manifests as API-only access. Model providers sell token-based inference rather than the weights themselves, with OpenAI generating \$5.7 billion in Q1~2026 revenue through token pricing ranging from \$0.05 to \$120 per million tokens~\cite{investingcom2026tokenpricing}. This mirrors Microsoft's historical practice of assessing per-device OEM licensing fees: the revenue is tied to usage rather than to the artifact itself.

DeepSeek's open-source R1 model, which demonstrated competitive benchmarks with OpenAI's o1, pioneered a ``third pathway'': generating revenue from open-source AI~\cite{thirdbridge2025deepseek}. Proprietary enterprise AI use reportedly shrank from 80\% to 44\% across 2025, while open-source alternatives gained ground~\cite{mckinsey2025aisurvey}. The \$200 billion-plus web of circular investments linking OpenAI, Anthropic, Microsoft, Google, NVIDIA, and Amazon through partnerships and equity stakes~\cite{brookings2025aicustomers} represents a vertically integrated \textit{keiretsu} structure that temporarily enforces the closed-source model through capital interdependence. It cannot, however, solve the fundamental copyability problem: weights, once leaked or distilled, propagate without marginal cost.

\subsection{Hardware Binding and Weight Protection: Rebuilding the Physical Moat}
\label{sec:hardware-binding}

If copyability undermines the software-only business model, the industry's response is to reconstruct a physical moat around model weights. Three strategies span a spectrum from software-based protection to permanent hardware binding.

The first approach treats weights as copyable but renders them useless without authorized hardware. RAND Corporation's 2024 playbook proposes 167 security measures for weight protection, recommending Trusted Execution Environment (TEE)-based encryption where weights are decrypted only inside a verified enclave~\cite{rand2024weightplaybook}. NVIDIA's H100, the first GPU with confidential computing, and Corvex's 2026 launch combining TEEs with hardware-encrypted GPU memory and Intel TDX~\cite{corvex2026secureweights}, exemplify this pathway: a form of digital rights management for AI models.

The second approach embeds hardware-level identity into the inference process. RePACK (Nature Communications, 2025) integrates Physical Unclonable Functions (PUFs) with ReRAM compute-in-memory, creating a device-specific fingerprint that ties weight execution to a particular piece of silicon~\cite{wu2025physicalunclonable}. This is an intermediate strategy: weights remain digitally distributed, but correct execution requires the matching hardware.

The third and most radical approach permanently burns model weights into silicon transistors. Taalas (2026) produces dedicated inference chips with a one-to-one chip-to-model binding, achieving approximately 17,000 tokens per second per user (roughly ten times GPU performance)~\cite{taalas_2026_hc1}. The model becomes inseparable from the hardware: the ``model is the computer.'' This creates a true physical moat analogous to Intel's chip fabrication, where the intellectual property is embodied in the manufacturing process itself. These three approaches (TEE encryption, PUF-based compute-in-memory, and permanent silicon burning) represent increasing degrees of ``physicalization,'' each trading flexibility for IP protection, directly paralleling the semiconductor industry's progression from licensed IP (ARM) to manufactured chips (Intel) to dedicated ASICs.

\subsection{The OS Layer as Traffic Gateway: Platform Strategy in the AI Era}
\label{sec:os-gateway}

The hardware-binding strategies of the previous subsection address weight protection at the model layer. Yet the history of the PC industry suggests that the highest-value capture point lies not in hardware or in the model, but in the platform layer: the interface that mediates between users and underlying capabilities.

Microsoft's pivotal 1990 achievement illustrates the dynamic. By securing OEM pre-installation agreements with thirty major PC manufacturers for Windows~3.0~\cite{licendi2024mslicensing}, Microsoft created a flywheel: ubiquitous distribution fed a developer ecosystem, which drove consumer lock-in, which sustained licensing revenue that reached \$23--29 billion per year in recent times~\cite{strategyzer2024msbusinessmodel}. Windows was the traffic gateway between users and PC computing. The entity that controlled the gateway captured disproportionate value, even though Intel manufactured the hardware that made it all possible.

Current AI chatbot interfaces, including ChatGPT, Claude, and Gemini, function as de facto ``operating systems'' for AI interaction. They are the traffic gateway between users and model capabilities. Walmart, Google, and WeChat are explicitly competing to become ``AI-era traffic gateways''~\cite{ceibs2025aiagents}, confirming that the platform layer is recognized as the highest-value capture point in the AI stack. Google's Android strategy, an open-source OS monetized through embedded services (ads, Play Store, cloud)~\cite{wikipedia2024msopensource}, represents an alternative model that AI platform companies may adopt: open the interface to gain market share, then monetize through embedded commerce and data flows. As AI industry standards and protocols mature, enabling model interoperability and portability, the entity controlling the gateway rather than the model itself will likely capture disproportionate value, exactly as Microsoft captured more value than Intel in the PC era~\cite{ccia2025competition}.

\subsection{Convergence Scenarios: When Training Stabilizes and Brains Become Physical Products}
\label{sec:convergence}

The semiconductor industry's four-stage evolution provides a roadmap for AI industry maturation. The prerequisite for full disaggregation is the emergence of industry standards and protocols that enable interoperability across layers, just as standard interfaces such as the x86 ISA, the PCI bus, and USB enabled the PC ecosystem to disaggregate~\cite{microchipusa2024amdhistory}.

Several signals suggest the industry is approaching a ``training stabilization'' inflection point. Gartner projects \$2.59 trillion in AI spending for 2026 (a 47\% year-over-year increase) and predicts major consolidation in agentic AI~\cite{gartner2026spending}. McKinsey reports that only 33\% of organizations have scaled AI beyond pilots~\cite{mckinsey2025aisurvey}, a ``trough of disillusionment'' that historically precedes industry maturation. The PC industry went through a similar phase: the 1983--1985 hardware shakeout eliminated dozens of incompatible platforms before the standard IBM PC compatible architecture enabled explosive growth.

When model architectures stabilize, the ``selling brains as physical devices'' model becomes viable. Dedicated inference chips with burned-in weights, such as the Taalas HC1 achieving tenfold GPU performance~\cite{taalas_2026_hc1}, could commoditize the model layer while creating a new hardware product category: the AI appliance. This would parallel the transition from mainframe time-sharing to personal computers, where computing ceased to be a shared service and became a personal possession.

The convergence of three trends (hardware binding technologies maturing, inference costs dropping 280-fold, and custom silicon proliferating) creates the conditions for a ``fabless AI'' industry structure. Model companies would design ``brains'' (architectures and weights), foundries (TSMC, Samsung) would manufacture them as dedicated chips, and platform companies would distribute them through OS-level gateways. The most likely end-state is a three-layer industry structure: fabless model designers (the future Intels and AMDs of AI), AI foundries manufacturing dedicated inference and training chips, and AI OS platforms controlling user relationships through interface and gateway services. This mirrors the semiconductor and PC disaggregation that created today's Intel, TSMC, and Microsoft, the very firms whose history has served as this section's predictive lens.
